\documentclass[A4,twocolumn]{article}

%\usepackage[iso]{umlaute}
%\usepackage{german}
\usepackage{hyperref}
\usepackage{graphicx}
\setlength{\parindent}{0cm}
\setlength{\parskip}{1ex}
\setlength{\columnsep}{25pt}

\textwidth=17cm
\textheight=23cm
\setlength{\unitlength}{0.5cm}
\setlength{\parindent}{0.0cm}
\setlength{\parskip}{1ex}
\raggedbottom
\sloppy
%\addtolength{\evensidemargin}{-5cm}
\addtolength{\oddsidemargin}{-1.5cm}
\addtolength{\topmargin}{-2cm}

\sloppy

% Your name
\author{Trung Nguyen\\ Technische Universit\"at M\"unchen}

\title{Seminar Cloud Computing \\
       {\bf From Concept to Production: Deploying TinyML in Industry}
}

% Date of your talk
\date{November 2024}


\begin{document}

\maketitle

\begin{abstract}


\end{abstract}

% \section defines numbered parts of the paper with titles
% there also are \subsection and \subsubsection
\section{Introduction}

% Use labels to be able to refer to this position from somewhere else
\label{introduction}

\subsection{Context and Importance of TinyML}
\subsection{Scope and Objective}




% free floating figure using width of one column.
\begin{figure}
\centerline{
\includegraphics[width=0.9\columnwidth]{resource/TUM-Logo-102.png}
}
\caption{The caption explaining what can be seen in the image/figure.
Readers often read captions first if they do not have much time. Thus,
it is important to find a good short explanation.}
% A label to allow refering to this figure in the text.
\label{TUM}
\end{figure}




\section{TinyML Overview (1 page)}

\subsection{Definition and Key Concepts}
\subsection{Why TinyML Matters}




\begin{equation}
a^2 + b^2 = c^2
\label{Pythagoras}
\end{equation}

Again, refering to this equation is easy (see Eq.~\ref{Pythagoras}).
If you do not need numbering for equations, use the {\em displaymath}
environment:

\begin{displaymath}
x_{1,2} = \frac{-b \pm \sqrt{b^2-4ac}}{2a}\\
\end{displaymath}




\section{Use Cases of TinyML}

Enumerations using bullet points:

\begin{itemize}
	\item 	IoT and Smart Devices
	\item 	Environmental Monitoring
	\item 	Industrial Applications
	\item 	IoT and Smart Devices
	\item 	Edge AI and Autonomous Systems

\end{itemize}



\section{Techniques in TinyML (3-4 pages)}



\begin{quote}
	``I think there is a world market for maybe five computers.''
	(T.J. Watson, IBM, 1943)
\end{quote}

The rest of the work (especially all the regular text) must be
written/phrased by you. If you write about some results or fact
stated in another paper, you should refer to it.
The `Analytical Engine'' --- a mechanical calculation machine ---
created by Charles Babbage in the year 1838 was based on the decimal
system
% use \cite to refer to papers from seminarpaper.bib
% this file is processed by bibtex, and it automatically adds numbering
\cite{hussein_original_2024, paul_rethinking_2021, de_prado_robustifying_2020,ren_synergy_2021,roshan_adaptive_2021}.

\section{Challenges and Future of TinyML }




\section{Conclusion}
\label{conclusion}



% Put citations from bibtex into References section which were not
% explicity cited.
\nocite{hussein_original_2024,paul_rethinking_2021}


\bibliographystyle{plain}
% Literature sources are to be found in seminarpaper.bib
\bibliography{seminarpaper}
\end{document}
